% Section: Recursive Auditing Systems
\section{Recursive Auditing Systems}
\label{sec:recursive-auditing}

% TODO: Expand this section.
%
% Key topics:
% - Definition of recursive auditing
% - Invariant specification and monitoring
% - Audit trail construction
% - Continuous vs. checkpoint-triggered auditing
% - Integration with milestone synchronization
% - Tooling and implementation approaches

Recursive auditing is the continuous process of validating system state against a
predefined set of invariants across all levels of the recursive development stack.
Unlike static analysis or point-in-time code review, recursive auditing operates
as an ongoing monitoring layer that accompanies each phase of autonomous execution.

\subsection{Invariant Specification}

An \textit{invariant} in the Reflector framework is a formally specified property
that must hold throughout the development process. Invariants are classified into
three categories:

\begin{itemize}
  \item \textbf{Structural invariants}: Properties of the codebase structure, such as
    module boundaries, interface contracts, and dependency constraints.
  \item \textbf{Behavioral invariants}: Properties of system behavior, such as test
    coverage thresholds, performance bounds, and security properties.
  \item \textbf{Governance invariants}: Properties of the development process itself,
    such as milestone completion order, approval record completeness, and agent scope
    compliance.
\end{itemize}

\subsection{Audit Trail}

The recursive audit trail is an append-only log of all significant system state
transitions, validation results, governance decisions, and agent actions. The audit
trail serves as the authoritative record for post-hoc analysis of drift incidents
and supports the reconstruction of the causal chain leading to any failure mode.

\subsection{Continuous Monitoring}

Reflector's auditing system monitors invariants continuously during autonomous
execution, not only at milestone boundaries. This enables early detection of
developing drift patterns before they cross the threshold of milestone validation
failure—providing an early warning system that can trigger precautionary governance
checkpoints before full drift occurs.
